\documentclass[8pt]{article}
\usepackage[a4paper,margin=0.9cm]{geometry}
\usepackage{amsmath,amssymb,mathtools,booktabs,longtable,array,adjustbox,pdflscape,microtype,ragged2e}
\usepackage[T1]{fontenc}
\setlength{\parindent}{0pt}
\setlength{\parskip}{0.45em}
\setlength{\tabcolsep}{2.5pt}
\renewcommand{\arraystretch}{1.10}
\setlength{\LTleft}{0pt}
\setlength{\LTright}{0pt}
\allowdisplaybreaks
\begin{document}
\begin{landscape}
\section*{T3 group factors: EFT\_2\_after\_S1\_S2}
Active theory: SMEFT with the Weinberg operator $C_5$.  No T3 heavy field remains dynamical.
\[j=\frac{d-1}{2},\qquad C_2(d)=\frac{d^2-1}{4},\qquad T(d)=\frac{d(d^2-1)}{12}.\]
In every comparison table, $\mathcal{S}$ denotes the RGE structure multiplying the displayed group factor in $16\pi^2\beta$.
\section*{Notation key}
\begin{longtable}{@{}p{0.19\linewidth}p{0.73\linewidth}@{}}
\toprule
symbol & interaction / definition \\
\midrule
$g_1$ & $D_\mu\supset i g_1 Y B_\mu$ \\
$g_2$ & $D_\mu\supset i g_2 T^A W_\mu^A$ \\
$g_3$ & $D_\mu\supset i g_3 t^A G_\mu^A$ \\
$Y_u$ & $-\bar Q\,Y_u\,\widetilde H\,u+\mathrm{h.c.}$ \\
$Y_d$ & $-\bar Q\,Y_d\,H\,d+\mathrm{h.c.}$ \\
$Y_e$ & $-\bar L\,Y_e\,H\,e+\mathrm{h.c.}$ \\
$T$ & $\operatorname{Tr}\!\left(Y_e^\dagger Y_e+3Y_u^\dagger Y_u+3Y_d^\dagger Y_d\right)$ \\
$\lambda_1$ & $\frac12\lambda_1(H^\dagger H)^2$ \\
$C_5^{ij}$ & $(L_i^T C L_j)\,H H$ \\
\bottomrule
\end{longtable}
\section*{Stage field content}
\begin{longtable}{@{}lccc p{0.20\linewidth}p{0.20\linewidth}@{}}
\toprule
run & $S_1$ & $S_2$ & $F$ & integrated at this stage & active after matching \\
\midrule
T3-A-m4 & $(1,-2)$ & $(3,-1)$ & $(2,- \frac{3}{2})$ & S1, S2 & none \\
T3-A-m2 & $(1,-1)$ & $(3,0)$ & $(2,- \frac{1}{2})$ & S1, S2 & none \\
T3-A-p0 & $(1,0)$ & $(3,1)$ & $(2,\frac{1}{2})$ & S1, S2 & none \\
T3-B-m3 & $(2,- \frac{3}{2})$ & $(2,- \frac{1}{2})$ & $(1,-1)$ & S1, S2 & none \\
T3-B-m1 & $(2,- \frac{1}{2})$ & $(2,\frac{1}{2})$ & $(1,0)$ & S1, S2 & none \\
T3-B-p1 & $(2,\frac{1}{2})$ & $(2,\frac{3}{2})$ & $(1,1)$ & S1, S2 & none \\
T3-C-m3 & $(2,- \frac{3}{2})$ & $(2,- \frac{1}{2})$ & $(3,-1)$ & S1, S2 & none \\
T3-C-m1 & $(2,- \frac{1}{2})$ & $(2,\frac{1}{2})$ & $(3,0)$ & S1, S2 & none \\
T3-C-p1 & $(2,\frac{1}{2})$ & $(2,\frac{3}{2})$ & $(3,1)$ & S1, S2 & none \\
T3-D-m2 & $(3,-1)$ & $(1,0)$ & $(2,- \frac{1}{2})$ & S1, S2 & none \\
T3-D-p0 & $(3,0)$ & $(1,1)$ & $(2,\frac{1}{2})$ & S1, S2 & none \\
T3-D-p2 & $(3,1)$ & $(1,2)$ & $(2,\frac{3}{2})$ & S1, S2 & none \\
T3-E-m4 & $(3,-2)$ & $(3,-1)$ & $(2,- \frac{3}{2})$ & S1, S2 & none \\
T3-E-m2 & $(3,-1)$ & $(3,0)$ & $(2,- \frac{1}{2})$ & S1, S2 & none \\
T3-E-p0 & $(3,0)$ & $(3,1)$ & $(2,\frac{1}{2})$ & S1, S2 & none \\
T3-E-p2 & $(3,1)$ & $(3,2)$ & $(2,\frac{3}{2})$ & S1, S2 & none \\
\bottomrule
\end{longtable}
\section*{Decoupling statement}
All dependence on the original $S_1$, $S_2$ and $F$ representations is contained in the matched boundary value $C_5(\mu_{\mathrm{th}})$.  The one-loop running below the final threshold is common to all T3 classes.
\subsection*{$ \beta_{C_5} $}
\textit{Run block 1: T3-A-m4--T3-B-m1}
\begin{longtable}{@{}p{0.31\linewidth}ccccc@{}}
\toprule
$\mathcal{S}$ & $ \mathrm{T3-A-m4}\;(\alpha=-4) $ & $ \mathrm{T3-A-m2}\;(\alpha=-2) $ & $ \mathrm{T3-A-p0}\;(\alpha=0) $ & $ \mathrm{T3-B-m3}\;(\alpha=-3) $ & $ \mathrm{T3-B-m1}\;(\alpha=-1) $ \\
\midrule
$ \lambda_1\,C_5 $ & $ 2 $ & $ 2 $ & $ 2 $ & $ 2 $ & $ 2 $ \\
$ g_2^2\,C_5 $ & $ -3 $ & $ -3 $ & $ -3 $ & $ -3 $ & $ -3 $ \\
$ T\,C_5 $ & $ 2 $ & $ 2 $ & $ 2 $ & $ 2 $ & $ 2 $ \\
$ (Y_e^\dagger Y_e)^T C_5 $ & $ - \frac{3}{2} $ & $ - \frac{3}{2} $ & $ - \frac{3}{2} $ & $ - \frac{3}{2} $ & $ - \frac{3}{2} $ \\
$ C_5(Y_e^\dagger Y_e) $ & $ - \frac{3}{2} $ & $ - \frac{3}{2} $ & $ - \frac{3}{2} $ & $ - \frac{3}{2} $ & $ - \frac{3}{2} $ \\
\bottomrule
\end{longtable}
\textit{Run block 2: T3-B-p1--T3-D-m2}
\begin{longtable}{@{}p{0.31\linewidth}ccccc@{}}
\toprule
$\mathcal{S}$ & $ \mathrm{T3-B-p1}\;(\alpha=1) $ & $ \mathrm{T3-C-m3}\;(\alpha=-3) $ & $ \mathrm{T3-C-m1}\;(\alpha=-1) $ & $ \mathrm{T3-C-p1}\;(\alpha=1) $ & $ \mathrm{T3-D-m2}\;(\alpha=-2) $ \\
\midrule
$ \lambda_1\,C_5 $ & $ 2 $ & $ 2 $ & $ 2 $ & $ 2 $ & $ 2 $ \\
$ g_2^2\,C_5 $ & $ -3 $ & $ -3 $ & $ -3 $ & $ -3 $ & $ -3 $ \\
$ T\,C_5 $ & $ 2 $ & $ 2 $ & $ 2 $ & $ 2 $ & $ 2 $ \\
$ (Y_e^\dagger Y_e)^T C_5 $ & $ - \frac{3}{2} $ & $ - \frac{3}{2} $ & $ - \frac{3}{2} $ & $ - \frac{3}{2} $ & $ - \frac{3}{2} $ \\
$ C_5(Y_e^\dagger Y_e) $ & $ - \frac{3}{2} $ & $ - \frac{3}{2} $ & $ - \frac{3}{2} $ & $ - \frac{3}{2} $ & $ - \frac{3}{2} $ \\
\bottomrule
\end{longtable}
\textit{Run block 3: T3-D-p0--T3-E-p0}
\begin{longtable}{@{}p{0.31\linewidth}ccccc@{}}
\toprule
$\mathcal{S}$ & $ \mathrm{T3-D-p0}\;(\alpha=0) $ & $ \mathrm{T3-D-p2}\;(\alpha=2) $ & $ \mathrm{T3-E-m4}\;(\alpha=-4) $ & $ \mathrm{T3-E-m2}\;(\alpha=-2) $ & $ \mathrm{T3-E-p0}\;(\alpha=0) $ \\
\midrule
$ \lambda_1\,C_5 $ & $ 2 $ & $ 2 $ & $ 2 $ & $ 2 $ & $ 2 $ \\
$ g_2^2\,C_5 $ & $ -3 $ & $ -3 $ & $ -3 $ & $ -3 $ & $ -3 $ \\
$ T\,C_5 $ & $ 2 $ & $ 2 $ & $ 2 $ & $ 2 $ & $ 2 $ \\
$ (Y_e^\dagger Y_e)^T C_5 $ & $ - \frac{3}{2} $ & $ - \frac{3}{2} $ & $ - \frac{3}{2} $ & $ - \frac{3}{2} $ & $ - \frac{3}{2} $ \\
$ C_5(Y_e^\dagger Y_e) $ & $ - \frac{3}{2} $ & $ - \frac{3}{2} $ & $ - \frac{3}{2} $ & $ - \frac{3}{2} $ & $ - \frac{3}{2} $ \\
\bottomrule
\end{longtable}
\textit{Run block 4: T3-E-p2--T3-E-p2}
\begin{longtable}{@{}p{0.31\linewidth}c@{}}
\toprule
$\mathcal{S}$ & $ \mathrm{T3-E-p2}\;(\alpha=2) $ \\
\midrule
$ \lambda_1\,C_5 $ & $ 2 $ \\
$ g_2^2\,C_5 $ & $ -3 $ \\
$ T\,C_5 $ & $ 2 $ \\
$ (Y_e^\dagger Y_e)^T C_5 $ & $ - \frac{3}{2} $ \\
$ C_5(Y_e^\dagger Y_e) $ & $ - \frac{3}{2} $ \\
\bottomrule
\end{longtable}
\section*{Flavor form}
\begin{align}
16\pi^2\frac{dC_5}{d\ln\mu}&=(2\lambda_1-3g_2^2+2T)C_5\\
&\quad-\frac32\left[(Y_e^\dagger Y_e)^T C_5+C_5(Y_e^\dagger Y_e)\right].
\end{align}
\end{landscape}
\end{document}